% Part: first-order-logic
% Chapter: syntax-and-semantics
% Section: satisfaction

\documentclass[../../../include/open-logic-section]{subfiles}

\begin{document}

\olfileid{fol}{syn}{sat}

\olsection{\article{structure} \printtoken{S}{structure}లో
  \article{formula} \printtoken{S}{formula} యొక్క సంతృప్తి}

\begin{explain}
ఒక వైపు పదాలు, \tetoken{సూత్రాలు}{!!{formula}s} వంటి వ్యక్తీకరణలను, మరో వైపు
\tetoken{నిర్మాణాలను}{!!{structure}s} అనుసంధానించే మౌలిక భావనలు: ఒక పదం యొక్క
\emph{\tetoken{విలువ}{!!{value}}}, \tetoken{ఒక సూత్రం}{!!a{formula}} యొక్క \emph{సంతృప్తి}. అనౌపచారికంగా,
ఒక పదం యొక్క \tetoken{విలువ}{!!{value}} అనేది \tetoken{ఒక నిర్మాణంలోని}{!!a{structure}} \tetoken{మూలకం}{!!{element}}---ఆ పదం
కేవలం ఒక స్థిరసంకేతమైతే, దానికి \tetoken{నిర్మాణం}{!!{structure}} నిర్దేశించిన వస్తువే దాని
\tetoken{విలువ}{!!{value}}; అది \tetoken{ప్రమేయ సంకేతాలను}{!!{function}s} ఉపయోగించి నిర్మించబడితే, ఆ పదంలోని
స్థిరసంకేతాల \tetoken{విలువలు}{!!{value}s}, ప్రమేయ సంకేతాలకు నిర్దేశించిన ప్రమేయాల నుంచి
దాని \tetoken{విలువను}{!!{value}} లెక్కిస్తాం. విధేయాలకు ఇచ్చిన అర్థనిర్దేశం ఆ
\tetoken{సూత్రాన్ని}{!!{formula}} \tetoken{నిర్మాణం}{!!{structure}} యొక్క వ్యక్తి క్షేత్రంలో సత్యం చేస్తే,
\tetoken{ఒక సూత్రం}{!!^a{formula}} ఆ \tetoken{ఒక నిర్మాణంలో}{!!a{structure}} \emph{సంతృప్తమవుతుంది}. ఈ సంతృప్తి
భావనను ఆగమనాత్మకంగా నిర్దేశిస్తాం: పరమాణు \tetoken{సూత్రాలు}{!!{formula}s} ఎప్పుడు
సంతృప్తమవుతాయో \tetoken{నిర్మాణం}{!!{structure}} యొక్క నిర్దేశమే నేరుగా చెబుతుంది; ఒక
సంక్లిష్ట \tetoken{సూత్రం}{!!{formula}} ఎప్పుడు సంతృప్తమవుతుందో, దాని ప్రధాన సంయోజకం లేదా
పరిమాణకం ఏదో, దాని తక్షణ \tetoken{ఉపసూత్రాలు}{!!{subformula}s} సంతృప్తమవుతాయో లేదో అనేదాని
ఆధారంగా నిర్వచిస్తాం.

ఇక్కడ పరిమాణకాల సందర్భం కాస్త సూక్ష్మమైనది. పరిమాణీకృత \tetoken{సూత్రం}{!!{formula}}
యొక్క తక్షణ \tetoken{ఉపసూత్రంలో}{!!{subformula}} ఒక స్వేచ్ఛా \tetoken{చరం}{!!{variable}} ఉంటుంది; కానీ
\tetoken{నిర్మాణాలు}{!!{structure}s}, \tetoken{చరాలు}{!!{variable}s} యొక్క \tetoken{విలువలను}{!!{value}s} నిర్దేశించవు. ఈ
కష్టాన్ని పరిష్కరించడానికి \emph{చర నిర్దేశాలను} కూడా ప్రవేశపెట్టి,
సంతృప్తిని కేవలం \tetoken{ఒక నిర్మాణానికి}{!!a{structure}} సాపేక్షంగా కాకుండా, \tetoken{ఒక నిర్మాణంతో}{!!a{structure}}
పాటు \tetoken{ఒక చరం}{!!a{variable}} నిర్దేశానికి సాపేక్షంగా నిర్వచిస్తాం.
\end{explain}

\begin{defn}[చర నిర్దేశం]
\tetoken{ఒక నిర్మాణం}{!!a{structure}}~$\Struct{M}$ కోసం \emph{చర నిర్దేశం}~$s$ అనేది ప్రతి
\tetoken{చరాన్ని}{!!{variable}} $\Domain M$లోని \tetoken{ఒక మూలకానికి}{!!a{element}} పంపే ప్రమేయం; అంటే,
$s\colon \Var \to \Domain M$.
\end{defn}

\begin{explain}
\tetoken{ఒక నిర్మాణం}{!!^a{structure}} ప్రతి \tetoken{స్థిరసంకేతానికి}{!!{constant}} \tetoken{ఒక విలువను}{!!a{value}}, చర నిర్దేశం ప్రతి
చరానికి ఒక విలువను నిర్దేశిస్తాయి. అయితే వాటి నుంచి నిర్మించిన పదాలు కూడా
\tetoken{వ్యక్తి క్షేత్రంలోని}{!!{domain}} \tetoken{మూలకాలను}{!!{element}s} సూచించాలని కోరుకుంటాం. ఇందుకోసం పదాల
\tetoken{విలువను}{!!{value}} ఆగమనాత్మకంగా నిర్వచిస్తాం. \tetoken{స్థిరసంకేతాలు}{!!{constant}s}, చరాల విషయంలో
విలువను వరుసగా \tetoken{నిర్మాణం}{!!{structure}}, చర నిర్దేశం చెప్పినట్లే తీసుకుంటాం; మరింత
సంక్లిష్టమైన పదాల విషయంలో, \tetoken{నిర్మాణం}{!!{structure}} \tetoken{ప్రమేయ సంకేతాలకు}{!!{function}s} నిర్దేశించిన
ప్రమేయాలను ఉపయోగించి పునరావర్తనంగా లెక్కిస్తాం.
\end{explain}

\begin{defn}[పదాల \tetoken{విలువ}{!!^{value}}]
$t$ అనేది భాష~$\Lang L$లోని పదం, $\Struct M$ అనేది $\Lang L$కు ఒక
\tetoken{నిర్మాణం}{!!{structure}}, $s$ అనేది $\Struct M$ కోసం \tetoken{ఒక చరం}{!!a{variable}} నిర్దేశం అయితే,
\emph{\tetoken{విలువ}{!!{value}}}~$\Value{t}{M}[s]$ను కింది విధంగా నిర్వచిస్తాం:
\begin{enumerate}
\item \indcase{t}{c}{$\Value{\indfrm}{M}[s] = \Assign{\indcomplex}{M}$.}
\item \indcase{t}{x}{$\Value{\indfrm}{M}[s] = s(\indcomplex)$.}
\item \indcase{t}{\Atom{f}{t_1, \ldots, t_n}}{
\[
\Value{\indfrm}{M}[s] = \Assign{f}{M}(\Value{t_1}{M}[s], \ldots,
\Value{t_n}{M}[s]).
\]}
\end{enumerate}
\end{defn}

\begin{defn}[$x$-భేదరూపం]
$s$ అనేది \tetoken{ఒక నిర్మాణం}{!!a{structure}}~$\Struct M$ కోసం \tetoken{ఒక చరం}{!!a{variable}} నిర్దేశం అయితే,
$x$కు ఏమి నిర్దేశిస్తుందనే విషయంలో మాత్రమే గరిష్ఠంగా $s$తో భేదించే
$\Struct M$ కోసం ఏ \tetoken{చరం}{!!{variable}} నిర్దేశం~$s'$నైనా, $s$ యొక్క
\emph{$x$-భేదరూపం} అంటాం. $s'$ అనేది $s$ యొక్క $x$-భేదరూపమైతే
$\varAssign{s'}{s}{x}$ అని రాస్తాం.
\end{defn}

\begin{explain}
ఒక నిర్దేశం~$s$ యొక్క $x$-భేదరూపం, $x$కు తప్పనిసరిగా \emph{వేరే}
దానిని నిర్దేశించాల్సిన అవసరం లేదని గమనించండి. వాస్తవానికి, ప్రతి
నిర్దేశం తనకే ఒక $x$-భేదరూపంగా పరిగణించబడుతుంది.
\end{explain}

\begin{defn}
  $s$ అనేది \tetoken{ఒక నిర్మాణం}{!!a{structure}}~$\Struct M$ కోసం \tetoken{ఒక చరం}{!!a{variable}} నిర్దేశం,
  $m \in \Domain{M}$ అయితే, $\Subst{s}{m}{x}$ అనే నిర్దేశం కింది విధంగా
  నిర్వచించిన చర నిర్దేశం:
  \[\Subst{s}{m}{x}(y) = \begin{cases}
    m & y \ident x \text{ అయితే}\\
    s(y) & \text{లేకపోతే}.
  \end{cases}\]
\end{defn}

మరోలా చెప్పాలంటే, $\Subst{s}{m}{x}$ అనేది $x$కు వ్యక్తి క్షేత్రపు
\tetoken{మూలకం}{!!{element}}~$m$ను నిర్దేశించే $s$ యొక్క ప్రత్యేక $x$-భేదరూపం; $x$
కాని \tetoken{చరాలకు}{!!{variable}s} మాత్రం $s$ నిర్దేశించిన వాటినే నిర్దేశిస్తుంది.

\begin{defn}[సంతృప్తి]
\ollabel{defn:satisfaction}
\tetoken{ఒక నిర్మాణం}{!!a{structure}}~$\Struct M$లో, \tetoken{ఒక చరం}{!!a{variable}} నిర్దేశం~$s$కు సాపేక్షంగా
\tetoken{ఒక సూత్రం}{!!a{formula}}~$!A$ యొక్క సంతృప్తిని $\Sat{M}{!A}[s]$ అనే సంకేతాలతో
కింది విధంగా పునరావర్తనంగా నిర్వచిస్తాం. (``$\Sat{M}{!A}[s]$ కాదు'' అని
అర్థం చేసుకోవడానికి $\Sat/{M}{!A}[s]$ అని రాస్తాం.)
\begin{enumerate}
\tagitem{prvFalse}{%
  \indcase{!A}{\lfalse}{$\Sat/{M}{\indfrm}[s]$.}}{}

\tagitem{prvTrue}{%
  \indcase{!A}{\ltrue}{$\Sat{M}{\indfrm}[s]$.}}{}

\item \indcase{!A}{\Atom{R}{t_1, \dots, t_n}}{$\Sat{M}{\indfrm}[s]$
  అవుతుంది అప్పుడే, అప్పుడు మాత్రమే $\langle \Value{t_1}{M}[s], \dots,
  \Value{t_n}{M}[s] \rangle \in \Assign{R}{M}$.}

\item \indcase{!A}{\eq[t_1][t_2]}{$\Sat{M}{\indfrm}[s]$ అవుతుంది
  అప్పుడే, అప్పుడు మాత్రమే $\Value{t_1}{M}[s] = \Value{t_2}{M}[s]$.}

\tagitem{prvNot}{%
  \indcase{!A}{\lnot !B}{$\Sat{M}{\indfrm}[s]$ అవుతుంది అప్పుడే,
    అప్పుడు మాత్రమే $\Sat/{M}{!B}[s]$.}}{}

\tagitem{prvAnd}{%
  \indcase{!A}{(!B \land !C)}{$\Sat{M}{\indfrm}[s]$ అవుతుంది అప్పుడే,
    అప్పుడు మాత్రమే $\Sat{M}{!B}[s]$ మరియు $\Sat{M}{!C}[s]$.}}{}

\tagitem{prvOr}{%
  \indcase{!A}{(!B \lor !C)}{$\Sat{M}{\indfrm}[s]$ అవుతుంది అప్పుడే,
    అప్పుడు మాత్రమే $\Sat{M}{!B}[s]$ లేదా $\Sat{M}{!C}[s]$ (లేదా రెండూ).}}{}

\tagitem{prvIf}{%
  \indcase{!A}{(!B \lif !C)}{$\Sat{M}{\indfrm}[s]$ అవుతుంది అప్పుడే,
    అప్పుడు మాత్రమే $\Sat/{M}{!B}[s]$ లేదా $\Sat{M}{!C}[s]$ (లేదా రెండూ).}}{}

\tagitem{prvIff}{%
  \indcase{!A}{(!B \liff !C)}{$\Sat{M}{\indfrm}[s]$ అవుతుంది అప్పుడే,
    అప్పుడు మాత్రమే $\Sat{M}{!B}[s]$, $\Sat{M}{!C}[s]$ రెండూ అవుతాయి,
    లేదా $\Sat{M}{!B}[s]$, $\Sat{M}{!C}[s]$ రెండూ కావు.}}{}

\tagitem{prvAll}{%
  \indcase{!A}{\lforall[x][!B]}{$\Sat{M}{\indfrm}[s]$ అవుతుంది అప్పుడే,
    అప్పుడు మాత్రమే ప్రతి \tetoken{మూలకం}{!!{element}}~$m \in \Domain M$కు
    $\Sat{M}{!B}[\Subst{s}{m}{x}]$.}}{}

\tagitem{prvEx}{%
  \indcase{!A}{\lexists[x][!B]}{$\Sat{M}{\indfrm}[s]$ అవుతుంది అప్పుడే,
  అప్పుడు మాత్రమే కనీసం ఒక \tetoken{మూలకం}{!!{element}}~$m \in \Domain M$కు
  $\Sat{M}{!B}[\Subst{s}{m}{x}]$.}}{}
\end{enumerate}
\end{defn}

\begin{explain}
చివరి \iftag{notprvEx,notprvAll}{నిబంధనలో}{రెండు నిబంధనల్లో} చర
నిర్దేశాలు ముఖ్యమైనవి.\iftag{prvAll}{ $\lforall[x][!B(x)]$ యొక్క
సంతృప్తిని ``అన్ని $m \in \Domain{M}$లకు $\Sat{M}{!B(m)}$'' అని
నిర్వచించలేం.}{}\iftag{prvEx}{ $\lexists[x][!B(x)]$ యొక్క సంతృప్తిని
``కనీసం ఒక $m \in \Domain{M}$కు $\Sat{M}{!B(m)}$'' అని
నిర్వచించలేం.}{} కారణం: $m \in \Domain M$ అయితే అది భాషలోని సంకేతం
కాదు; కాబట్టి $!B(m)$ ఒక \tetoken{సూత్రం}{!!a{formula}} కాదు (అంటే,
$\Subst{!B}{m}{x}$ నిర్వచించబడలేదు). అంతేకాక, $\Struct{M}$లోని ప్రతి
\tetoken{మూలకానికి}{!!{element}} పేరు పెట్టే \tetoken{స్థిరసంకేతాలు}{!!{constant}s} లేదా పదాలు మనకు అందుబాటులో
ఉంటాయని భావించలేం; \tetoken{నిర్మాణాలు}{!!{structure}s} నిర్వచనంలో అలాంటి నిబంధనేమీ లేదు.
ప్రామాణిక భాషలో \tetoken{స్థిరసంకేతాలు}{!!{constant}s} సమితి \tetoken{అనంతంగా లెక్కించదగిన}{!!{denumerable}}; అందువల్ల
$\Domain{M}$ \tetoken{లెక్కించదగిన}{!!{enumerable}} కాకపోతే, ప్రతి వస్తువుకు పేరు పెట్టడానికి
సరిపడా \tetoken{స్థిరసంకేతాలు}{!!{constant}s} కూడా ఉండవు.

చరాలను వ్యక్తి క్షేత్రంలోని \tetoken{మూలకాలతో}{!!{element}s} నేరుగా అనుసంధానించే
\tetoken{చరం}{!!{variable}} నిర్దేశాలను ప్రవేశపెట్టి ఈ సమస్యను పరిష్కరిస్తాం. అప్పుడు,
ఉదాహరణకు, \iftag{prvEx}{$\lexists[x][!B(x)]$}{$\lforall[x][!B(x)]$}
అనేది~$\Struct M$లో సంతృప్తమవుతుంది అప్పుడే, అప్పుడు మాత్రమే
\iftag{prvEx}{కనీసం ఒక $m \in \Domain{M}$కు}{అన్ని $m \in
\Domain{M}$లకు, $\Sat{M}{!B(m)}$} అని చెప్పడానికి బదులుగా, అది
$s$కు \emph{సాపేక్షంగా}~$\Struct M$లో సంతృప్తమవుతుంది అప్పుడే, అప్పుడు
మాత్రమే, $!B(x)$ అనేది \iftag{prvEx}{కనీసం ఒక}{ప్రతి}
$m \in \Domain M$కు $\Subst{s}{m}{x}$కు సాపేక్షంగా సంతృప్తమవుతుంది
అని చెబుతాం.
\end{explain}

\begin{ex}
$\Lang{L} = \{a, b, f, R\}$ అనుకుందాం; ఇక్కడ $a$, ~$b$లు
\tetoken{స్థిరసంకేతాలు}{!!{constant}s}, $f$ ఒక ద్విస్థాన \tetoken{ప్రమేయ సంకేతం}{!!{function}}, $R$ ఒక ద్విస్థాన
\tetoken{విధేయ సంకేతం}{!!{predicate}}. కింది విధంగా నిర్వచించిన \tetoken{నిర్మాణం}{!!{structure}}~$\Struct{M}$ను
పరిశీలించండి:
\begin{enumerate}
\item $\Domain M = \{1, 2, 3, 4\}$
\item $\Assign{a}{M} = 1$
\item $\Assign{b}{M} = 2$
\item $\Assign{f}{M}(x, y) = x+y$, $x+y \le 3$ అయితే; లేకపోతే $= 3$.
\item $\Assign{R}{M} = \{\tuple{1, 1}, \tuple{1, 2}, \tuple{2, 3}, \tuple{2, 4}\}$
\end{enumerate}
ప్రతి \tetoken{చరానికి}{!!{variable}} $1 \in \Domain{M}$ను నిర్దేశించే ప్రమేయం
$s(x) = 1$, $\Struct{M}$ కోసం ఒక చర నిర్దేశం.

అప్పుడు
\begin{align*}
\Value{f(a,b)}{M}[s] & = \Assign{f}{M}(\Value{a}{M}[s], \Value{b}{M}[s]).
\intertext{$a$, ~$b$లు \tetoken{స్థిరసంకేతాలు}{!!{constant}s} కాబట్టి, $\Value{a}{M}[s]
  = \Assign{a}{M} = 1$, అలాగే $\Value{b}{M}[s] = \Assign{b}{M} = 2$. కనుక}
\Value{f(a,b)}{M}[s] & = \Assign{f}{M}(1, 2) = 1+2 = 3.
\intertext{$f(f(a,b),a)$ విలువను లెక్కించడానికి మనం పరిశీలించాల్సింది}
\Value{f(f(a,b),a)}{M}[s] & = \Assign{f}{M}(\Value{f(a, b)}{M}[s],
\Value{a}{M}[s]) = \Assign{f}{M}(3, 1) = 3,
\intertext{ఎందుకంటే $3+1 > 3$. $s(x) = 1$, అలాగే $\Value{x}{M}[s] =
  s(x)$ కాబట్టి, మనకు ఇది కూడా లభిస్తుంది:}
\Value{f(f(a,b),x)}{M}[s] & = \Assign{f}{M}(\Value{f(a, b)}{M}[s],
\Value{x}{M}[s]) = \Assign{f}{M}(3, 1) = 3,
\end{align*}

ఒక పరమాణు \tetoken{సూత్రం}{!!{formula}}~$R(t_1, t_2)$, దాని ఆర్గ్యుమెంట్ల విలువల క్రమిత
జంట, అంటే $\tuple{\Value{t_1}{M}[s], \Value{t_2}{M}[s]}$,
$\Assign{R}{M}$లో \tetoken{ఒక మూలకం}{!!a{element}} అయితే సంతృప్తమవుతుంది. కాబట్టి,
ఉదాహరణకు, $\tuple{\Value{b}{M}, \Value{f(a,b)}{M}} = \tuple{2, 3}
\in \Assign{R}{M}$ కావడంతో $\Sat{M}{R(b,f(a,b))}[s]$; కానీ
$\tuple{1, 3} \notin \Assign{R}{M}$ కావడంతో
$\Sat/{M}{R(x, f(a,b))}[s]$.\sourcecorrection{OLTEFOLSAT-001}{మూలంలో సంబంధ అర్థనిర్దేశం $\Assign{R}{M}$కు పొరపాటున జోడించిన $[s]$ను తొలగించాం; చర నిర్దేశం సూత్ర సంతృప్తికే వర్తిస్తుంది.}

పరమాణువుకాని సూత్రం~$!A$ సంతృప్తమవుతుందో నిర్ధారించడానికి, దాని
ప్రధాన సంయోజకానికి వర్తించే ఆగమనాత్మక నిర్వచన నిబంధనను ఉపయోగిస్తారు.
ఉదాహరణకు, $R(a, a) \lif (R(b, x) \lor R(x, b))$లో ప్రధాన సంయోజకం
$\lif$; అందువల్ల
\begin{align*}
  & \Sat{M}{R(a, a) \lif (R(b, x) \lor R(x, b))}[s] \text{ అవుతుంది అప్పుడే, అప్పుడు మాత్రమే }\\
  & \qquad
  \Sat/{M}{R(a,a)}[s] \text{ లేదా } \Sat{M}{R(b, x) \lor R(x, b)}[s]
  \intertext{$\Sat{M}{R(a,a)}[s]$ కాబట్టి ($\tuple{1,1} \in
    \Assign{R}{M}$ కావడం వల్ల), సమాధానాన్ని ఇంకా నిర్ణయించలేం;
    మొదట $\Sat{M}{R(b, x) \lor R(x, b)}[s]$ అవుతుందో తెలుసుకోవాలి:}
  & \Sat{M}{R(b, x) \lor R(x, b)}[s] \text{ అవుతుంది అప్పుడే, అప్పుడు మాత్రమే }\\
  & \qquad \Sat{M}{R(b,x)}[s] \text{ లేదా } \Sat{M}{R(x, b)}[s]
  \intertext{$\Sat{M}{R(x, b)}[s]$ కాబట్టి ఇది జరుగుతుంది
    ($\tuple{1,2} \in \Assign{R}{M}$ కావడం వల్ల).}
\end{align*}

$s$ యొక్క $x$-భేదరూపం అంటే, $x$కు ఏమి నిర్దేశిస్తుందనే విషయంలో మాత్రమే
గరిష్ఠంగా $s$తో భేదించే చర నిర్దేశమని గుర్తుచేసుకోండి. $\Domain{M}$లోని
ప్రతి \tetoken{మూలకానికి}{!!{element}} $s$ యొక్క ఒక $x$-భేదరూపం ఉంది:
\begin{align*}
  s_1 & = \Subst{s}{1}{x}, &
  s_2 & = \Subst{s}{2}{x},\\
  s_3 & = \Subst{s}{3}{x}, &
  s_4 & = \Subst{s}{4}{x}.
\end{align*}
ఉదాహరణకు, $s_2(x) = 2$; అలాగే $x$ కాని అన్ని చరాలు~$y$కు
$s_2(y) = s(y) = 1$. $\Domain{M} = \{1, 2, 3, 4\}$ కాబట్టి,
$\Struct{M}$ నిర్మాణం కోసం $s$కు ఉన్న $x$-భేదరూపాలు ఇవన్నీ. ముఖ్యంగా,
$s_1 = s$ అని గమనించండి ($s$ ఎల్లప్పుడూ తనకే ఒక $x$-భేదరూపం).

\iftag{prvEx}{అస్తిత్వ పరిమాణీకృత \tetoken{సూత్రం}{!!{formula}}~$\lexists[x][!A(x)]$
  సంతృప్తమవుతుందో నిర్ణయించడానికి, కనీసం ఒక $m \in \Domain M$కు
  $\Sat{M}{!A(x)}[\Subst{s}{m}{x}]$ అవుతుందో నిర్ణయించాలి. కాబట్టి,
  \[
  \Sat{M}{\lexists[x][(R(b,x) \lor R(x,b))]}[s],
  \]
  ఎందుకంటే $\Sat{M}{R(b,x) \lor R(x, b)}[\Subst{s}{1}{x}]$
  ($\Subst{s}{3}{x}$ కూడా ఈ అవసరాన్ని తీరుస్తుంది). కానీ,
  \[
  \Sat/{M}{\lexists[x][(R(b,x) \land R(x,b))]}[s]
  \]
  ఎందుకంటే, మనం ఏ $m \in \Domain{M}$ను ఎంచుకున్నా,
  $\Sat/{M}{R(b,x) \land R(x,b)}[\Subst{s}{m}{x}]$.}{}

\iftag{prvAll}{సార్వత్రిక పరిమాణీకృత \tetoken{సూత్రం}{!!{formula}}~$\lforall[x][!A(x)]$
  సంతృప్తమవుతుందో నిర్ణయించడానికి, అన్ని $m \in \Domain M$లకు
  $\Sat{M}{!A(x)}[\Subst{s}{m}{x}]$ అవుతుందో నిర్ణయించాలి. కాబట్టి,
  \[
  \Sat{M}{\lforall[x][(R(x,a) \lif R(a,x))]}[s],
  \]
  ఎందుకంటే అన్ని $m \in \Domain M$లకు
  $\Sat{M}{R(x,a) \lif R(a,x)}[\Subst{s}{m}{x}]$. $m = 1$కు
  $\Sat{M}{R(a,x)}[\Subst{s}{1}{x}]$ అవుతుంది, కనుక పర్యవసానం సత్యం;
  $m = 2$, ~$3$, ~$4$లకు $\Sat/{M}{R(x,a)}[\Subst{s}{m}{x}]$
  అవుతుంది, కనుక పూర్వపక్షం అసత్యం. కానీ,
  \[
  \Sat/{M}{\lforall[x][(R(a,x) \lif R(x,a))]}[s]
  \]
  ఎందుకంటే $\Sat/{M}{R(a,x) \lif R(x,a)}[\Subst{s}{2}{x}]$
  ($\Sat{M}{R(a, x)}[\Subst{s}{2}{x}]$ మరియు $\Sat/{M}{R(x,
  a)}[\Subst{s}{2}{x}]$ కావడం వల్ల).}{}

\iftag{defEx}{అస్తిత్వ పరిమాణీకృత \tetoken{సూత్రం}{!!{formula}}~$\lexists[x][!A(x)]$
  సంతృప్తమవుతుందో నిర్ణయించడానికి,
  $\Sat{M}{\lnot \lforall[x][\lnot !A(x)]}[s]$ అవుతుందో నిర్ణయించాలి.
  ఉదాహరణకు, మనకు
  \[
  \Sat{M}{\lexists[x][(R(b,x) \lor R(x,b))]}[s].
  \]
  మొదట, $\Sat{M}{R(b,x) \lor R(x, b)}[\Subst{s}{1}{x}]$
  ($\Subst{s}{3}{x}$ కూడా ఈ అవసరాన్ని తీరుస్తుంది). కనుక,
  $\Sat/{M}{\lnot(R(b,x) \lor R(x,b))}[\Subst{s}{1}{x}]$; అందువల్ల
  $\Sat/{M}{\lforall[x][\lnot((R(b,x) \lor R(x,b))]}[s]$; కాబట్టి
  $\Sat{M}{\lnot\lforall[x][\lnot((R(b,x) \lor
  R(x,b))]}[s]$. మరోవైపు,
  \[
  \Sat/{M}{\lexists[x][(R(b,x) \land R(x,b))],}[s].
  \]
  కారణం: $\Sat{M}{\lforall[x][\lnot(R(b,x) \land
  R(x,b))]}[s]$, ఎందుకంటే ఏ $m \in \Domain M$కూ
  $\Sat{M}{R(b,x) \land R(x,b)}[\Subst{s}{m}{x}]$ కాదు. ఈ ఉదాహరణల
  నుంచి మీరు ఊహించగలిగినట్లుగా, కనీసం ఒక $m \in \Domain M$కు
  $\Sat{M}{!A(x)}[\Subst{s}{m}{x}]$ అవుతుంది అప్పుడే, అప్పుడు మాత్రమే
  $\Sat{M}{\lexists[x][!A(x)]}[s]$.}{}

\iftag{defAll}{సార్వత్రిక పరిమాణీకృత \tetoken{సూత్రం}{!!{formula}}~$\lforall[x][!A(x)]$
  సంతృప్తమవుతుందో నిర్ణయించడానికి,
  $\Sat{M}{\lnot\lexists[x][\lnot !A(x)]}[s]$ అవుతుందో నిర్ణయించాలి.
  ఉదాహరణకు,
  \[
  \Sat{M}{\lforall[x][(R(x,a) \lif R(a,x))]}[s],
  \]
  మొదట, అన్ని $m \in \Domain M$లకు
  $\Sat{M}{R(x,a) \lif R(a,x)}[\Subst{s}{m}{x}]$;
  ($\Sat{M}{R(a,x)}[\Subst{s}{1}{x}]$, అలాగే $m = 2$, ~$3$, ~$4$లకు
  $\Sat/{M}{R(x,a)}[\Subst{s}{m}{x}]$).\sourcecorrection{OLTEFOLSAT-002}{ఈ పర్యవసాన నిరూపణలో $m=2,3,4$కు అసత్యమయ్యేది పూర్వపక్షం $R(x,a)$; మూలంలో ఉన్న $R(a,x)$ను సరిచేశాం.}
  అందువల్ల $\Sat{M}{\lnot(R(x,a) \lif R(a,x))}[\Subst{s}{m}{x}]$
  అయ్యే $m \in \Domain M$ లేదు; కనుక
  $\Sat/{M}{\lexists[x][\lnot(R(x,a) \lif R(a,x))]}[s]$. అందుచేత,
  $\Sat{M}{\lnot\lexists[x][\lnot(R(x,a) \lif R(a,x))]}[s]$.
  మరోవైపు,
  \[
  \Sat/{M}{\lforall[x][(R(a,x) \lif R(x,a))]}[s],
  \]
  ఎందుకంటే $\Sat/{M}{R(a,x) \lif R(x,a)}[\Subst{s}{2}{x}]$; కాబట్టి
  $\Sat{M}{\lexists[x][\lnot(R(a,x) \lif R(x,a))]}[s]$. ఈ
  ఉదాహరణల నుంచి మీరు ఊహించగలిగినట్లుగా, ప్రతి $m \in \Domain M$కు
  $\Sat{M}{!A(x)}[\Subst{s}{m}{x}]$ అవుతుంది అప్పుడే, అప్పుడు మాత్రమే
  $\Sat{M}{\lforall[x][!A(x)]}[s]$.}{}

మరింత సంక్లిష్టమైన సందర్భంగా దీనిని పరిశీలించండి:
\[
\lforall[x][(R(a,x) \lif \lexists[y][R(x,y)])].
\]
$\Sat/{M}{R(a,x)}[\Subst{s}{3}{x}]$, అలాగే
$\Sat/{M}{R(a,x)}[\Subst{s}{4}{x}]$ కాబట్టి, షరతు యొక్క పర్యవసానం
గురించి పరిశీలించాల్సిన సందర్భాలు $m = 1$, ~$m = 2$ మాత్రమే.\sourcecorrection{OLTEFOLSAT-003}{మూలంలో రెండవ సందర్భం వద్ద తప్పిపోయిన చరాన్ని పునరుద్ధరించి ``$m=2$''గా చదివాం.}
$\Sat{M}{\lexists[y][R(x,y)]}[\Subst{s}{1}{x}]$ అవుతుందా? కనీసం ఒక
$n \in \Domain M$కు
$\Sat{M}{R(x,y)}[\Subst{\Subst{s}{1}{x}}{n}{y}]$ అయితే అవుతుంది.
వాస్తవానికి $n = 1$ను తీసుకుంటే,
$\Subst{\Subst{s}{1}{x}}{n}{y} = \Subst{s}{1}{y} = s$. $s(x) = 1$,
$s(y) = 1$, ఇంకా $\tuple{1,1} \in \Assign{R}{M}$ కాబట్టి సమాధానం అవును.

$\Sat{M}{\lexists[y][R(x,y)]}[\Subst{s}{2}{x}]$ అవుతుందో
నిర్ణయించడానికి, $\Subst{\Subst{s}{2}{x}}{n}{y}$ అనే \tetoken{చరం}{!!{variable}}
నిర్దేశాలను పరిశీలించాలి. ఇక్కడ $n = 1$కు ఈ నిర్దేశం
$s_2 = \Subst{s}{2}{x}$; ఇది $R(x,y)$ను సంతృప్తిపరచదు ($s_2(x) = 2$,
$s_2(y) = 1$, ఇంకా $\tuple{2,1}\notin \Assign{R}{M}$). అయితే,
$\Subst{\Subst{s}{2}{x}}{3}{y} = \Subst{s_2}{3}{y}$ను పరిశీలించండి.
$\tuple{2,3} \in \Assign{R}{M}$ కాబట్టి
$\Sat{M}{R(x,y)}[\Subst{s_2}{3}{y}]$; అందువల్ల
$\Sat{M}{\lexists[y][R(x,y)]}[s_2]$.

కాబట్టి ప్రతి $m \in \Domain M$కు,
$m = 3$, ~$4$ అయితే $\Sat/{M}{R(a,x)}[\Subst{s}{m}{x}]$, లేదా
$m = 1$, ~$2$ అయితే $\Sat{M}{\lexists[y][R(x,y)]}[\Subst{s}{m}{x}]$;
అందువల్ల\sourcecorrection{OLTEFOLSAT-004}{బాహ్య సార్వత్రిక పరిమాణకం $x$కు సంబంధించిన ఈ సంగ్రహంలో మూలంలోని ``for all $n$''ను ``for all $m$''గా సరిచేశాం.}
\[
\Sat{M}{\lforall[x][(R(a,x) \lif \lexists[y][R(x,y)])]}[s].
\]
మరోవైపు,
\[
\Sat/{M}{\lexists[x][(R(a,x) \land \lforall[y][R(x,y)])]}[s].
\]
$m = 1$, ~$m = 2$కు మాత్రమే $\Sat{M}{R(a,x)}[\Subst{s}{m}{x}]$.
కానీ $m$ యొక్క ఈ రెండు విలువల్లో ప్రతి దానికీ, తిరిగి ఒక
$n \in \Domain M$ ఉంది---అదే $n = 4$---దానికి
$\Sat/{M}{R(x,y)}[\Subst{\Subst{s}{m}{x}}{n}{y}]$; అందువల్ల
$m = 1$, ~$m = 2$కు
$\Sat/{M}{\lforall[y][R(x,y)]}[\Subst{s}{m}{x}]$. మొత్తానికి,
$\Sat{M}{R(a,x) \land \lforall[y][R(x,y)]}[\Subst{s}{m}{x}]$ అయ్యే
$m \in \Domain M$ ఏదీ లేదు.
\end{ex}

\iftag{defEx}{%
\begin{prop}\ollabel{prop:sat-ex}
$\Sat{M}{\lexists[x][!B(x)]}[s]$ అవుతుంది అప్పుడే, అప్పుడు మాత్రమే $s$కు
$\Sat{M}{!B(x)}[s']$ అయ్యే ఒక $x$-భేదరూపం $s'$ ఉంది.
\end{prop}

\begin{proof}
  వ్యాయామం.
\end{proof}
}{}

\tagprob{defEx}
\begin{prob}
  \olref[fol][syn][sat]{prop:sat-ex}ను నిరూపించండి.
\end{prob}
\tagendprob

\iftag{defAll}{%
\begin{prop}\ollabel{prop:sat-all}
$\Sat{M}{\lforall[x][!B(x)]}[s]$ అవుతుంది అప్పుడే, అప్పుడు మాత్రమే $s$
యొక్క ప్రతి $x$-భేదరూపం~$s'$కు $\Sat{M}{!B(x)}[s']$.
\end{prop}

\begin{proof}
  వ్యాయామం.
\end{proof}
}{}

\tagprob{defAll}
\begin{prob}
  \olref[fol][syn][sat]{prop:sat-all}ను నిరూపించండి.
\end{prob}
\tagendprob

\begin{prob}
$\Lang L = \{c, f, A\}$లో ఒక \tetoken{స్థిరసంకేతం}{!!{constant}}, ఒక ఏకస్థాన
\tetoken{ప్రమేయ సంకేతం}{!!{function}}, ఒక ద్విస్థాన \tetoken{విధేయ సంకేతం}{!!{predicate}} ఉన్నాయని, \tetoken{నిర్మాణం}{!!{structure}}
$\Struct{M}$ను కింది విధంగా ఇచ్చారని అనుకుందాం:
\begin{enumerate}
\item $\Domain M = \{1, 2, 3\}$
\item $\Assign{c}{M} = 3$
\item $\Assign{f}{M}(1) = 2, \Assign{f}{M}(2) = 3, \Assign{f}{M}(3) = 2$
\item $\Assign{A}{M} = \{\tuple{1, 2}, \tuple{2, 3}, \tuple{3, 3}\}$
\end{enumerate}
(a) అన్ని \tetoken{చరాలు}{!!{variable}s}~$v$కు $s(v) = 1$ అనుకుందాం. కిందిది
అవుతుందో లేదో కనుగొనండి:
\[
\Sat{M}{\lexists[x][(A(f(z), c) \lif \lforall[y][(A(y, x) \lor A(f(y),
      x))])]}[s]
\]
ఎందుకు అవుతుందో లేదా ఎందుకు కాదో వివరించండి.

(b) ఆ \tetoken{సూత్రం}{!!{formula}} సంతృప్తమవని మరొక నిర్మాణాన్ని, \tetoken{చరం}{!!{variable}}
నిర్దేశాన్ని ఇవ్వండి.
\end{prob}

\end{document}
